% mono/preamble/environments.tex
%
% Segment environments — the visual register for the 19 segment types
% defined in FORMAT.md. Architecture:
%
%   \segmenthead{<Type Label>}{<Title>}{<status>}
%      ... segment body ...
%   \segmentfoot
%
% The kramdown emitter produces these directly. No 19 distinct LaTeX
% environments — one parameterized macro pair, type label as a string. This
% keeps the AAD type taxonomy and the LaTeX taxonomy decoupled: when FORMAT
% grows a 20th type, the build pipeline maps it to the right label and the
% rendering needs no change here.
%
% Numbering uses a single shared counter per chapter (matches the existing
% bin/build markdown behavior: "Definition I.3", "Result II.7", etc.). The
% counter is named 'segment' rather than 'claim' (per FORMAT's vocabulary —
% segments hold claims; not all segments contain a single claim).

\newcounter{segment}[chapter]
% FORMAT.md uses "Definition I.3, Result II.7" with roman chapter numerals;
% match it here. (kaobook's kao chapter style already renders the BIG margin
% chapter glyph in roman, but \thechapter itself is arabic — \Roman{chapter}
% in the counter format aligns the body convention with the head.)
\renewcommand{\thesegment}{\Roman{chapter}.\arabic{segment}}

% cleveref format for our custom `segment` counter. Without this,
% cleveref doesn't know how to render \cref{seg:...} and falls back to
% `??` — even though the label resolves and the counter is correct.
%
% IMPORTANT: cleveref's format placeholder for the stored label value
% is `#1`, NOT `\thesegment`. Using `\thesegment` would expand at cref-
% TIME (giving the current counter value) instead of label-TIME (the
% value that was recorded when \label was emitted). Every cref would
% then render the same most-recent counter — sending all clicks to the
% wrong segment regardless of which label they targeted. `#2#1#3` is
% the canonical "hyperref-open + value + hyperref-close" idiom.
%
% Format prints just the counter value (`I.4`, `II.7`, …) with no type
% prefix. Authors include the type word in source — "see Definition
% #def-foo" — and the cref expands to the number alone, avoiding the
% doubling that a fixed prefix would cause across 19 segment types.
\AtBeginDocument{%
  \crefformat{segment}{#2#1#3}%
  \Crefformat{segment}{#2#1#3}%
  \crefrangeformat{segment}{#3#1#4--#5#2#6}%
  \Crefrangeformat{segment}{#3#1#4--#5#2#6}%
  \crefmultiformat{segment}{#2#1#3}{, #2#1#3}{, #2#1#3}{, #2#1#3}%
  \Crefmultiformat{segment}{#2#1#3}{, #2#1#3}{, #2#1#3}{, #2#1#3}%
}

% Internal section subheads — Formal Expression / Epistemic Status /
% Discussion / Findings. Each renders as a full-width thin rule across
% the segment band, with the subhead label right-aligned beneath. The
% rule does the visual section break; the right-aligned label names it
% quietly without competing with body text for the reader's eye.
%
% Orphan policy: reserve vertical space for the rule, the label, and a
% few lines of the section that follows; discourage a break in between.
\NewDocumentCommand{\segmentsubhead}{m}{%
  \needspace{5\baselineskip}%
  \par\medskip%
  \noindent\textcolor{rule-gray}{\rule{\segmentrulewidth}{0.3pt}}\par%
  \nobreak\nointerlineskip%
  \vskip-0.4\baselineskip%
  \noindent\makebox[\segmentrulewidth][r]{\sffamily\bfseries #1}\par%
  \nobreak\smallskip%
  \nopagebreak[4]%
}

% Width that spans body + margin gutter + margin column. kaobook's body
% column is \textwidth and the margin column sits to its right with
% \marginparsep gap and \marginparwidth width. We compute the union so
% segment-header rules read as page-spanning visual separators rather than
% body-column decorations.
\newcommand{\segmentrulewidth}{\dimexpr\textwidth+\marginparsep+\marginparwidth\relax}

% Segment opener. Three args:
%   #1 = type label (string, e.g. "Result", "Derivation")
%   #2 = title (the human-readable form of the slug)
%   #3 = status (string, e.g. "exact"; empty string = no badge)
%
% The third argument is always required; the kramdown emitter passes the
% empty string when frontmatter has no `status:` set, which suppresses the
% badge via the IfStrEq below.
%
% Segment-header palette — eight pastel tints across the 19 segment
% types. Updated from a more saturated first cut to a subtler pastel set;
% hex values estimated from a reference image and rounded to neat values.
\definecolor{tint-1}{HTML}{DDE2D2}  % pale sage
\definecolor{tint-2}{HTML}{D2D7DC}  % pale gray-blue
\definecolor{tint-3}{HTML}{DCD9D7}  % warm pale gray
\definecolor{tint-4}{HTML}{D5E5D7}  % pale mint
\definecolor{tint-5}{HTML}{CFE0EC}  % pale sky
\definecolor{tint-6}{HTML}{F2EAD0}  % pale cream
\definecolor{tint-7}{HTML}{F4DECF}  % pale peach
\definecolor{tint-8}{HTML}{F1D5DA}  % pale pink

% Live alias the segmenthead box reads from; \setheadertint{Label} updates
% it before each segment is drawn so each tcolorbox sees its own color.
\colorlet{header-tint}{tint-1}

% Dispatch by type-label (the human-readable form passed to \segmenthead's
% first argument): Definition, Result, Derivation, etc. Order follows the
% palette mapping Joseph chose for the live corpus distribution; the two
% currently-unused FORMAT types (Corollary, Aside) take sensible defaults
% (Corollary→Result family, Aside→Detail family) so they Just Work when
% they land in source.
% Mid-gray used for the bullet between type-label and title in segment
% headers. Light enough to feel like punctuation, dark enough to be
% clearly visible (the rule-gray we previously used was effectively
% invisible against the tinted backgrounds).
\definecolor{bullet-mid}{HTML}{666666}

% Stage glyph — promotion-stage indicator on the far right of the segment
% header strip in review mode. Five filled/empty dots map to the five
% post-`old` promotion stages from FORMAT.md. `missing`/`old` yield no
% glyph (those segments don't reach \segmenthead anyway — they go
% through \missingsegment).
% \CIRCLE (filled) and \Circle (empty) from wasysym are size-paired and
% reliably available — both render at the same baseline size, unlike
% \textbullet/\textcircled{} or the EB-Garamond-missing Unicode ●/○.
\NewDocumentCommand{\stageglyph}{m}{%
  {\color{bullet-mid}\footnotesize\IfStrEqCase{#1}{%
    {draft}           {\CIRCLE\,\Circle\,\Circle\,\Circle\,\Circle}%
    {deps-verified}   {\CIRCLE\,\CIRCLE\,\Circle\,\Circle\,\Circle}%
    {claims-verified} {\CIRCLE\,\CIRCLE\,\CIRCLE\,\Circle\,\Circle}%
    {format-clean}    {\CIRCLE\,\CIRCLE\,\CIRCLE\,\CIRCLE\,\Circle}%
    {candidate}       {\CIRCLE\,\CIRCLE\,\CIRCLE\,\CIRCLE\,\CIRCLE}%
  }[]}%
}

\NewDocumentCommand{\setheadertint}{m}{%
  \IfStrEqCase{#1}{%
    {Definition}      {\colorlet{header-tint}{tint-1}}%
    {Formulation}     {\colorlet{header-tint}{tint-2}}%
    {Detail}          {\colorlet{header-tint}{tint-2}}%
    {Discussion}      {\colorlet{header-tint}{tint-3}}%
    {Measurement}     {\colorlet{header-tint}{tint-3}}%
    {Derived}         {\colorlet{header-tint}{tint-4}}%
    {Derivation}      {\colorlet{header-tint}{tint-4}}%
    {Postulate}       {\colorlet{header-tint}{tint-4}}%
    {Scope}           {\colorlet{header-tint}{tint-5}}%
    {Normative}       {\colorlet{header-tint}{tint-5}}%
    {Hypothesis}      {\colorlet{header-tint}{tint-6}}%
    {Result}          {\colorlet{header-tint}{tint-7}}%
    {Observation}     {\colorlet{header-tint}{tint-7}}%
    {Worked Example}  {\colorlet{header-tint}{tint-7}}%
    {Empirical}       {\colorlet{header-tint}{tint-7}}%
    {Corollary}       {\colorlet{header-tint}{tint-7}}%  unused; share Result
    {Sketch}          {\colorlet{header-tint}{tint-8}}%
    {Proposed Schema} {\colorlet{header-tint}{tint-8}}%
    {Aside}           {\colorlet{header-tint}{tint-2}}%  unused; share Detail
  }[\colorlet{header-tint}{tint-1}]%
}

% Title is set large + bold (upright) so it reads with weight comparable
% to inline body-bold cues like "Lemma 1." or "Proof." — competing with
% the prose rather than looking decorative.
%
% Upper rule is 2.5x thicker and ~black, lower rule stays thin and gray —
% the asymmetry pulls the eye to the segment opening and lets the closing
% rule fade to a quiet separator. Between the rules: a tcolorbox holding
% the type-label + status-badge + title line on a soft tinted background.
%
% Orphan/widow policy: \needspace asks for enough vertical space to hold
% the header + a few lines of the first paragraph; if it's not available,
% LaTeX breaks the page first and the segment opens on the next page with
% its body intact. \nopagebreak after the bottom rule discourages a break
% between header and the paragraph that follows.
% Fourth arg is the segment's `stage:` frontmatter value, used in review
% mode for the dot-progression glyph on the far right of the header strip.
% Empty string elides it (public-variant builds, or stage-less segments).
\NewDocumentCommand{\segmenthead}{m m m m}{%
  \needspace{10\baselineskip}%
  \refstepcounter{segment}%
  \setheadertint{#1}%
  \par\medskip%
  \noindent\begin{tcolorbox}[
    enhanced,
    width            = \segmentrulewidth,
    colback          = header-tint,
    boxrule          = 0pt,
    arc              = 0pt,
    left             = 8pt, right  = 8pt,
    top              = 6pt, bottom = 6pt,
    before skip      = 0pt, after skip = 0pt,
    % Both rules black; top 2x the bottom so the eye registers the
    % segment-opening edge over the closing edge while both stay clearly
    % visible against the tinted background.
    borderline north = {1.5pt}{0pt}{black},
    borderline south = {0.75pt}{0pt}{black},
  ]
  {\sffamily\bfseries #1~\thesegment}%
  \IfStrEq{#3}{}{}{\statusbadge{#3}}%
  \nobreak\quad\textcolor{bullet-mid}{\textbullet}\quad%
  {\large\bfseries #2}.%
  \IfStrEq{#4}{}{}{\hfill\stageglyph{#4}}%
  \end{tcolorbox}\par%
  \nobreak\medskip%
  \nopagebreak[4]%
}

% Segment closer — paragraph spacing only; no visible rule. The next
% segment's opener provides its own top border, so a foot rule would only
% produce a doubled separator when segments sit back-to-back. A chapter
% break or page break is enough of a closer for the trailing-segment case.
\NewDocumentCommand{\segmentfoot}{}{%
  \par\medskip%
}

% (Stage display moved into the segment header strip as a dot glyph;
% see \stageglyph and \segmenthead's fourth argument. This command kept
% only for backward compatibility in case any caller still emits it.)
\NewDocumentCommand{\segmentstage}{m}{}

% Missing-segment placeholder — for slugs listed in OUTLINE.md but without
% a corresponding src/ file. Generates a TOC entry and a visible marker so
% reviewers see the gap in the rendered output, not just in a sidecar list.
% Rules are the same full-width treatment as \segmenthead so the visual
% break is consistent for landed and missing segments.
\NewDocumentCommand{\missingsegment}{m m m}{%
  % #1 = type label, #2 = slug, #3 = claim text (from the outline table)
  \needspace{8\baselineskip}%
  \refstepcounter{segment}%
  % Emit the seg: label even though the segment isn't written — cross-refs
  % from other segments still resolve to a number rather than ??.
  \label{seg:#2}%
  \par\medskip%
  \noindent{\color{status-tentative}\rule{\segmentrulewidth}{0.75pt}}\par%
  \nobreak\smallskip%
  \noindent%
  {\sffamily\bfseries\color{status-tentative} #1~\thesegment}%
  \quad\textcolor{bullet-mid}{\textbullet}\quad%
  {\large\bfseries\color{status-tentative} #3}.%
  \par\nobreak\smallskip%
  \noindent\emph{Segment \texttt{#2} not yet written.}\par%
  \nobreak\smallskip%
  \noindent{\color{status-tentative}\rule{\segmentrulewidth}{0.75pt}}\par%
  \nobreak\medskip%
}

% Gap marker — for --GAP-- entries in OUTLINE tables (theory not yet
% developed; no slug, just a description of the open question).
\NewDocumentCommand{\outlinegap}{m}{%
  \par\medskip%
  \noindent\textcolor{status-tentative}{\textbf{[Gap]}}~\emph{#1}\par%
  \medskip%
}

% Segment epigraph — summary block between the segment title and the
% first \segmentsubhead ("Formal Expression"). Slightly smaller, justified,
% centered horizontally on the full segment band (\segmentrulewidth).
%
% Implementation: capture the body into an lrbox, then place it inside a
% \makebox of width \segmentrulewidth with center-alignment. The \makebox
% is rendered \noindent so it begins flush with the body column's left
% edge; its visual extent then runs from x=0 to x=\segmentrulewidth, with
% the inner parbox (width 0.7 * segmentrulewidth) horizontally centered.
%
% Why not adjustwidth*: on chapter-first pages the kao chapter style
% transitions through \pagelayout{wide}, leaving the page in a state
% where adjustwidth* samples wider dimensions than the surrounding body
% prose actually uses, and the epigraph then extends past the page's
% left edge. \makebox doesn't sample those page-state dimensions; the
% positioning is fully determined by the explicit width we hand it.
\newsavebox{\segepibox}

\NewDocumentEnvironment{segmentepigraph}{}{%
  % \medskip leading (was \smallskip) so the gap above the epigraph
  % matches the gap above the header and the gap below the epigraph
  % (header→epigraph and epigraph→first-subhead now both read as
  % section-level breaks rather than tightly-stacked content).
  \par\medskip%
  \begin{lrbox}{\segepibox}%
  \begin{minipage}{0.7\segmentrulewidth}%
  \small%
  % A bit of parskip so multi-paragraph epigraphs breathe — minipage
  % defaults to \parskip=0 which is too tight when the summary runs to
  % two or three paragraphs.
  \setlength{\parskip}{0.8\baselineskip}%
  \setlength{\parindent}{0pt}%
  \noindent\ignorespaces%
}{%
  \par%
  \end{minipage}%
  \end{lrbox}%
  \noindent\makebox[\segmentrulewidth][c]{\usebox{\segepibox}}%
  \par\medskip%
}

% Wide-section wrapper — for Discussion and Findings, where the section
% content extends to the full segment band (body + margin gutter +
% margin column). No background, no rule; just a wider text area.
% tcolorbox's `blanker` skin gives us width control without decoration.
%
% Parskip is bumped inside the wrapper so multi-paragraph commentary in
% Discussion / Findings (which is interpretive prose, often several
% paragraphs of explanation) reads with proper breathing room between
% paragraphs rather than the body-default tight stacking.
\NewDocumentEnvironment{segmentwidesection}{}{%
  \par\medskip%
  \begin{tcolorbox}[
    blanker, breakable,
    width        = \segmentrulewidth,
    left         = 0pt, right  = 0pt,
    top          = 0pt, bottom = 0pt,
    before skip  = 0pt, after skip = 0pt,
  ]%
  \setlength{\parskip}{0.8\baselineskip}%
  \setlength{\parindent}{0pt}%
}{%
  \end{tcolorbox}%
  \par\medskip%
}

% Callouts — Obsidian-style `> [!type] title` blocks. We tint by type
% (warning/note/info/tip/danger), set a thin colored left border, and
% render the title in bold. Body content (paragraphs, lists, tables) gets
% the small/dim register so the callout reads as a contextual note rather
% than competing with body prose.
\definecolor{callout-warning-bg}    {HTML}{FDF5D9}
\definecolor{callout-warning-rule}  {HTML}{C49A2E}
\definecolor{callout-note-bg}       {HTML}{E7F0FB}
\definecolor{callout-note-rule}     {HTML}{3D6FB5}
\definecolor{callout-tip-bg}        {HTML}{E6F4EA}
\definecolor{callout-tip-rule}      {HTML}{3A8050}
\definecolor{callout-danger-bg}     {HTML}{FBE6E6}
\definecolor{callout-danger-rule}   {HTML}{B53A3A}

\NewDocumentCommand{\calloutpalette}{m}{%
  \IfStrEqCase{#1}{%
    {warning}{\colorlet{callout-bg}{callout-warning-bg}\colorlet{callout-rule}{callout-warning-rule}}%
    {note}   {\colorlet{callout-bg}{callout-note-bg}   \colorlet{callout-rule}{callout-note-rule}}%
    {info}   {\colorlet{callout-bg}{callout-note-bg}   \colorlet{callout-rule}{callout-note-rule}}%
    {tip}    {\colorlet{callout-bg}{callout-tip-bg}    \colorlet{callout-rule}{callout-tip-rule}}%
    {danger} {\colorlet{callout-bg}{callout-danger-bg} \colorlet{callout-rule}{callout-danger-rule}}%
  }[\colorlet{callout-bg}{callout-note-bg}\colorlet{callout-rule}{callout-note-rule}]%
}

% \begin{callout}{type}{title}  ...body...  \end{callout}
\NewDocumentEnvironment{callout}{m m}{%
  \calloutpalette{#1}%
  \par\medskip%
  \begin{tcolorbox}[
    enhanced, breakable,
    width        = \segmentrulewidth,
    colback      = callout-bg,
    colframe     = callout-bg,
    colupper     = bullet-mid,
    fontupper    = \small,
    boxrule      = 0pt,
    arc          = 0pt,
    left         = 12pt, right  = 12pt,
    top          = 8pt,  bottom = 8pt,
    before skip  = 0pt,  after skip = 0pt,
    borderline west = {2.5pt}{0pt}{callout-rule},
  ]%
  {\sffamily\bfseries\color{callout-rule}\MakeUppercase{#1}%
   \IfStrEq{#2}{}{}{\quad\normalcolor #2}}\par\smallskip%
}{%
  \end{tcolorbox}%
  \par\medskip%
}

% Working Notes — review-mode-only environment for the trailing ## Working
% Notes section of each segment. Background-tinted box spanning body +
% margin gutter + margin column (so the notes visibly step out of the main
% column register), shrunk and dimmed so the eye reads them as process
% chatter rather than load-bearing prose. The renderer only opens the env
% if the segment actually has a Working Notes heading, so no empty boxes.
\usepackage{tcolorbox}
\tcbuselibrary{breakable}

\definecolor{wn-bg}{HTML}{FAF7E5}  % subtle cream-tinted off-white
\definecolor{wn-fg}{HTML}{6E6354}  % warm muted gray-brown

% colupper / fontupper carry the text styling as persistent properties of
% the tcolorbox rather than as inline `\color` / `\small` directives —
% which is what lets them survive a page break inside a breakable box.
% (The inline form resets to black when the second tcolorbox fragment
% begins on the next page, a recurring KOMA + tcolorbox surprise.)
\NewDocumentEnvironment{workingnotes}{}{%
  \par\bigskip%
  \begin{tcolorbox}[
    enhanced, breakable,
    width        = \segmentrulewidth,
    colback      = wn-bg,
    colframe     = wn-bg,
    colupper     = wn-fg,
    fontupper    = \small,
    boxrule      = 0pt,
    arc          = 0pt,
    left         = 12pt, right  = 12pt,
    top          = 10pt, bottom = 10pt,
    before skip  = 0pt,  after skip = 0pt,
  ]%
  {\sffamily\bfseries Working Notes.}\par\smallskip%
}{%
  \end{tcolorbox}%
  \par\medskip%
}
